\begin{lemma}\label{prop:T-comparison}
	考虑来自伴随对 $(F, G)$ 的单子 $T$ 和余单子 $L$. 设 $N \in \Obj(\mathcal{D})$, 则资料
	\[ M := G(N) \in \Obj(\mathcal{C}), \quad a: T(M) = GFG(N) \xrightarrow{G\varepsilon_N} G(N) = M \]
	给出 $(M, a) \in \Obj(\mathcal{C}^T)$. 由此得到函子 $\mathbb{K}: \mathcal{D} \to \mathcal{C}^T$, 满足 $G = U^T \mathbb{K}$.
	
	对偶地, $F$ 也有典范分解 $\mathcal{C} \xrightarrow{\mathbb{K}} \mathcal{D}^L \to \mathcal{D}$, 今后写作 $F = U^L \mathbb{K}$, 其中 $U^L$ 是忘却函子\footnote{采用相同符号 $\mathbb{K}$ 不至于引起太大的混淆, 因为本书不会同时考虑单子和余单子.}.
\end{lemma}
\begin{proof}
	只论 $T$ 的情形. 定义 \ref{def:T-Alg} 的图表对此化为
	\[\begin{tikzcd}
		GFGFG(N) \arrow[r, "GFG\varepsilon" inner sep=0.5em] \arrow[d, "{G\varepsilon FG}"'] & GFG(N) \arrow[d, "G\varepsilon"] \\
		GFG(N) \arrow[r, "G\varepsilon"'] & G(N)
	\end{tikzcd} \quad \begin{tikzcd}
		G(N) \arrow[r, "\eta G"] \arrow[equal, rd] & GFG(N) \arrow[d, "{G\varepsilon}"] \\
		& G(N),
	\end{tikzcd} \]
	交换性的验证和例 \ref{eg:adjunction-monad} 如出一辙, 不必重复; $N \mapsto (M, a)$ 的函子性是自明的.
\end{proof}

对于给定伴随对 $(F, G)$, 应用函子是一个丢失结构的过程. 引理 \ref{prop:T-comparison} 将 $G$ (或 $F$) 拆成 $U^T \mathbb{K}$ (或 $U^L \mathbb{K}$). 我们想明白结构是否只被第二步的忘却函子 $U^T$ 丢失, 如果答案是肯定的, 过程中丢失的信息便能借助单子 $T$ (或余单子 $L$) 来重构. 这就启发了以下概念.

\begin{definition}\label{def:monadic}
	考虑一对伴随函子
	$\begin{tikzcd}
		F: \mathcal{C} \arrow[shift left, r] & \mathcal{D}: G \arrow[shift left, l]
	\end{tikzcd}$,
	以此定义 $\mathcal{C}$ 上的单子 $T$ 和 $\mathcal{D}$ 上的余单子 $L$.
	\begin{itemize}
		\item 若引理 \ref{prop:T-comparison} 的函子 $\mathbb{K}: \mathcal{D} \to \mathcal{C}^T$ 是范畴等价, 则称此伴随对是\emph{单子的}.
		\item 若对偶版本 $\mathbb{K}: \mathcal{C} \to \mathcal{D}^L$ 是范畴等价, 则称此伴随对是\emph{余单子}的.
	\end{itemize}
\end{definition}
